\subsection{\highlight{Computational Efficiency}}

\subsubsection{Token Consumption}

We further analyze the computational cost of different baselines by measuring token consumption across input, output and total tokens. The results are reported in Table \ref{table:token_consumption}. Values are reported as mean token consumption over all problem instances. Standard deviation is omitted due to aggregation-level logging constraints.

\begin{table}[h]
\centering
\small
\begin{tabular}{lccc}
\hline
\textbf{Baseline} & \textbf{Input tokens} & \textbf{Output tokens} & \textbf{Total tokens} \\
\hline
Single-agent & 20 & 2130 & 2151 \\
Chain-of-thought & 1444 & 2904 & 4349 \\
Actor-Critic & 4572 & 6865 & 11438 \\
ReAct & 292244 & 11876 & 304121 \\
Majority Voting & 5201 & 8527 & 13728 \\
Debate & 54285 & 27717 & 82003 \\
FP2MP & 111323 & 29150 & 140474 \\
\hline
\end{tabular}
\caption{Token consumption across baselines (mean values over 60 problem instances).}
\label{table:token_consumption}
\end{table}

The computational analysis reveals substantial differences in cost across baselines.
Single-Agent and Chain-of-Thought exhibit the lowest token consumption due to their direct generation structure, while iterative and multi-agent approaches entail significantly higher computational overhead.

Among the evaluated methods, ReAct demonstrates the highest token consumption, primarily due to repeated reasoning-action-observation cycles and extensive interaction with external tools.
Debate also require substantial computational resources since multiple agents iteratively refine and exchange intermediate solutions.

FP2MP exhibits higher computational cost compared to simpler baselines, reflecting the complexity of blackboard-based coordination and tool calling mechanism.
However, this additional cost is accompanied by improved performance across reflective reasoning dimensions, suggesting a trade-off between computational efficiency and reasoning quality in ill-defined problem solving.

\subsubsection{\highlight{API and Tool Calls}}

We report the mean number of API and external tool calls across all baselines and FP2MP variants in Table \ref{table:api_tool_calls}.

\begin{table}[h]
\centering
\small
\begin{tabular}{lcc}
\hline
\textbf{Baseline} & \textbf{API calls}& \textbf{Tool calls}\\
\hline
Single-agent & 1.0 ± 0.0& 0.0 ± 0.0\\
Chain-of-thought & 2.0 ± 0.0& 0.0 ± 0.0\\
Actor-Critic & 3.0 ± 0.0& 0.0 ± 0.0\\
ReAct & 18.88 ± 16.21& 32.63 ± 17.64\\
Majority Voting & 4.0 ± 0.0& 0.0 ± 0.0\\
Debate & 13.0 ± 0.0& 0.0 ± 0.0\\
FP2MP & 66.07 ± 13.21& 42.82 ± 11.82\\
\hline
FP2MP (w/o blackboard) & 99.7 ± 10.86& 67.38 ± 9.83\\
FP2MP (w/o critic) & 62.83 ± 11.24& 41.07 ± 10.16\\
FP2MP (w/o ReDI) & 42.32 ± 14.56& 27.15 ± 11.09\\
FP2MP (w/o tools) & 13.55 ± 1.85& 0.0 ± 0.0\\
\hline
\end{tabular}
\caption{\highlight{Number of API and external tool calls} across baselines and FP2MP variants (mean ± std over 60 problem instances).}
\label{table:api_tool_calls}
\end{table}

The results show that simpler baselines require substantially fewer interactions, while ReAct incurs additional overhead due to its iterative tool-augmented reasoning process.
FP2MP requires the largest number of API and tool calls among the evaluated methods, reflecting its multi-agent coordination and iterative tool use.
The ablation results further show that removing the blackboard increases the number of interactions, suggesting that centralized coordination can reduce the number of agent interactions.
In contrast, removing ReDI or tool access substantially reduces the number of calls, while removing the critic has a comparatively small effect on overall interaction volume.